\subsection*{Debug RTTI and \texttt{debugger}}

When \texttt{print} / \texttt{println} is not enough, Croqtile supports debug RTTI (Runtime Type Information) that makes its types visible to \texttt{debugger}.

Compile with \texttt{-g} to enable debug symbols:

\begin{lstlisting}[language={},caption={Debug RTTI and debugger}]
croqtile kernel.co -g -o kernel_debug
\end{lstlisting}

The generated code includes RTTI structs for Croqtile types:

\begin{table}[htbp]
\centering\small
\begin{tabularx}{\linewidth}{@{}lXX@{}}
\toprule
\textbf{Croqtile type} & \textbf{GDB type} & \textbf{Fields} \\
\midrule
Shaped tensor \textless{}br/\textgreater{} (\texttt{s32 [M, N]}) & \texttt{croq:\allowbreak{}:\allowbreak{}rtti:\allowbreak{}:\allowbreak{}spanned\textless{}int, 2\textgreater{}} & \texttt{.span.data[]} (dimensions), \textless{}br/\textgreater{}\texttt{.stride.data[]} (strides), \textless{}br/\textgreater{}\texttt{.data} (pointer) \\
Index tuple & \texttt{croq:\allowbreak{}:\allowbreak{}rtti:\allowbreak{}:\allowbreak{}bounded\_\allowbreak{}ituple\textless{}N\textgreater{}} & \texttt{.data[]} (values), \textless{}br/\textgreater{}\texttt{.ub[]} (upper bounds) \\
Integer tuple & \texttt{croq:\allowbreak{}:\allowbreak{}rtti:\allowbreak{}:\allowbreak{}ituple\textless{}N\textgreater{}} & \texttt{.data[]} (values) \\
Multi-dim span & \texttt{croq:\allowbreak{}:\allowbreak{}rtti:\allowbreak{}:\allowbreak{}mdspan\textless{}N\textgreater{}} & \texttt{.data[]} (extents) \\
\bottomrule
\end{tabularx}
\end{table}

\begin{lstlisting}[language={},caption={Debug RTTI and debugger}]
debugger -q ./kernel_debug
(gdb) break my_kernel
(gdb) run
(gdb) ptype __dbg_lhs
type = struct croq::rtti::spanned<int, 2>
(gdb) print __dbg_lhs.span.data[0]
$1 = 32
(gdb) print __dbg_lhs.span.data[1]
$2 = 64
(gdb) print __dbg_lhs.data != 0
$3 = true
\end{lstlisting}

The \texttt{\_\_dbg\_} prefix on variable names is generated by the compiler --- it makes Croqtile variables visible alongside the generated C++ intermediates.
